\section{Related Work}
\label{sec:related}
\subsection{LEO Network Measurements}
The first measurement studies of Starlink reported median RTTs of a few tens of milliseconds with frequent spikes and loss under load~\cite{michel2022first}. Later studies traced many of the latency jumps to a global scheduler that reassigns users to satellites every 15\,s~\cite{mohan2024multifaceted,tanveer2023making}. Long campaigns characterized the access, point-of-presence, and backbone segments~\cite{pan2023measuring}, and Starlink throughput was shown to vary on time scales from sub-second to diurnal~\cite{garcia2023multitimescale}. OneWeb has no inter-satellite links in service, and its latency changes with every satellite and gateway handover~\cite{zhao2025measuring}. A study on the same testbed as ours measured OneWeb dynamics and the benefit of combining paths in a liquid data transport~\cite{perrin2025measurement}. These studies explain why LEO link state changes quickly but do not forecast it at the sub-second scale needed by a packet scheduler.

\subsection{Network Performance Prediction}
Learned throughput prediction is well established for cellular links. Lumos5G maps location, mobility, and radio context to 5G throughput~\cite{narayanan2020lumos5g}, and recurrent models predict mobile bandwidth one step ahead~\cite{mei2019realtime}. Classical and deep models have also been compared on cellular traces~\cite{raca2020leveraging}. A study of sub-second cellular throughput found that learned models gain little over simple predictors at fine time scales~\cite{basit2024predictability}. For LEO networks, recent work predicts Starlink throughput from months of measurements~\cite{liu2025vivisecting} and global Starlink latency and throughput at coarse resolution~\cite{benghe2026horizon}. LEO RTT has been predicted with spatiotemporal attention~\cite{tian2025attention}. A concurrent preprint forecasts LEO link state with physics-informed features and quantile bands~\cite{chi2026reading}. All of these report point accuracy on test data, and none checks the causality of the input alignment or calibrates uncertainty on unseen sessions. Our evaluation includes persistence, seasonal, and linear references, following the observation that simple linear forecasters are often competitive with deep sequence models~\cite{zeng2023transformers}.

\subsection{Sequence Models and Probabilistic Forecasting}
Our forecaster builds on the encoder-decoder framework~\cite{sutskever2014sequence} with gated recurrent units~\cite{cho2014learning}. It uses additive attention~\cite{bahdanau2015neural} and is trained with scheduled teacher forcing~\cite{bengio2015scheduled}. Forecast accuracy is reported relative to a naive forecast, as recommended for series of different scales~\cite{hyndman2006another}. For probabilistic forecasts, a proper scoring rule rewards honest uncertainty~\cite{gneiting2007strictly}, and the optimal point forecast depends on the loss used to evaluate it~\cite{gneiting2011making}. Quantile regression with the pinball loss estimates conditional quantiles directly~\cite{koenker1978regression}. Deep ensembles improve both accuracy and uncertainty~\cite{lakshminarayanan2017simple}.

\subsection{Conformal Prediction}
Split conformal prediction gives intervals with finite-sample coverage when calibration and test data are exchangeable~\cite{lei2018distribution}. CQR adapts the interval width to the input by calibrating quantile regression outputs~\cite{romano2019conformalized}. Conformal methods have been applied to multi-horizon recurrent forecasts~\cite{stankeviciute2021conformal} and extended to non-exchangeable data with explicit coverage gaps~\cite{barber2023conformal}. Adaptive versions track distribution shift online~\cite{gibbs2021adaptive}. We use CQR across measurement sessions and report the coverage achieved on separate sessions.

\subsection{Multipath Scheduling}
MPQUIC extends QUIC to several paths~\cite{deconinck2017multipath}. Its default scheduler sends on the path with the lowest smoothed RTT, a policy inherited from MPTCP~\cite{paasch2014experimental}. BLEST and ECF account for head-of-line blocking and completion time on heterogeneous paths~\cite{ferlin2016blest,lim2017ecf}. Peekaboo learns a scheduling policy online~\cite{wu2020peekaboo}. MAMS uses mobility and handover awareness~\cite{yang2024mams}, and path-state forecasts have been fed to a 5G multi-access scheduler~\cite{maglione2025intelligent}. These schedulers act on point estimates. The calibrated intervals in this paper provide the risk information that they lack.
